%====================================================================================================
\documentclass[fontsize=12pt]{article}
%====================================================================================================
\thispagestyle{empty}
\renewcommand{\baselinestretch}{1.2}
\setlength{\tabcolsep}{2pt}
%====================================================================================================
\usepackage[a4paper, top=10mm, left=7.5mm]{geometry}
\usepackage[colorlinks,bookmarks=false,linkcolor=black,urlcolor=black]{hyperref}
\usepackage{multirow}
\usepackage{graphicx}
\usepackage{lipsum}
%====================================================================================================
% Insert an empty line with a negative vertical space, reducing the spacing typically added by \\.
% /!\ Tune this number so that your content fits the page length
\newcommand{\emptyLine}{\\[-0.2cm]}
% Define a bullet point with horizontal space on both sides
\newcommand{\bulletpoint}{\hspace*{\labelsep}\textbullet\hspace*{\labelsep}}
% Create a simple entry spanning 2 columns, left-aligned
\newcommand{\simpleEntry}[1]{\multicolumn{2}{l}{#1}}
% Create a section title, enlarged text spanning 2 columns
\newcommand{\sectionTitle}[1]{\emptyLine\emptyLine\simpleEntry{\Large{#1}}}
% Create a title entry, bold text spanning 2 columns
\newcommand{\titleEntry}[1]{\emptyLine\simpleEntry{\textbf{#1}}}
% Create an indented entry spanning 2 columns
\newcommand{\indentedEntry}[1]{\simpleEntry{\hspace{0.2cm}{#1}}}
% Create an entry with a bullet point, indented, spanning 2 columns
\newcommand{\bulletpointEntry}[1]{\indentedEntry{\bulletpoint{#1}}}
% Create a multicolumn entry spanning 3 columns in a table, preceded by an empty line.
\newcommand{\profileEntry}[1]{\emptyLine\multicolumn{3}{p{\textwidth}}{#1}}
%====================================================================================================
\begin{document}
\begin{tabular}{p{5cm}lr}
%====================================================================================================
\simpleEntry{\Large{\textbf{John SMITH}}} & \multirow[c]{3}{*}{\fbox{\includegraphics[height=3cm]{picture}}} \\
\simpleEntry{\large{\emph{Profession title}}} \\
\emptyLine
Street & Age (Birth date) \\
ZIP City & Nationality \\
Country \\
012 345 678 \\
john.smith@email.com \\
%====================================================================================================
\sectionTitle{Profile} \\ \hline
 \profileEntry{
 \lipsum[1][1-4]
 } \\
%====================================================================================================
\sectionTitle{Education} \\ \hline
 \titleEntry{U2, University 2, City, Country} \\
  \indentedEntry{Name of the diploma} & $20xx-20xx$ \\
 \titleEntry{U1, University 1, City, Country} \\
  \indentedEntry{Name of the diploma} & $20xx-20xx$ \\
  \indentedEntry{Name of the diploma} & $20xx-20xx$ \\
%====================================================================================================
\sectionTitle{Work experience} \\ \hline 
 \titleEntry{C2, Company 2, City, Country} \\
  \indentedEntry{Name of the position} & $20xx-20xx$ \\
   \bulletpointEntry{\lipsum[1][1]} \\
   \bulletpointEntry{\lipsum[1][1]} \\
  \indentedEntry{Name of the position} & $20xx-20xx$ \\
   \bulletpointEntry{\lipsum[1][1]} \\
 \titleEntry{C1, Company 1, City, Country} \\
  \indentedEntry{Name of the position} & $20xx-20xx$ \\
   \bulletpointEntry{\lipsum[1][1]} \\
%====================================================================================================
\sectionTitle{Computer skills} \\ \hline 
 \emptyLine
 {\textbf{Skill 1}} & Skill \\
 {\textbf{Skill 2}} & Skill \\
 {\textbf{Skill 3}} & Skill \\
%====================================================================================================
\sectionTitle{Languages} \\ \hline
 \emptyLine
 {\textbf{Language 1}} & Level \\
 {\textbf{Language 2}} & Level \\
%====================================================================================================
\sectionTitle{Interests} \\ \hline 
 \titleEntry{Activity} \\
  \bulletpointEntry{\lipsum[1][1]} & $20xx-20xx$ \\
%====================================================================================================
\end{tabular}
\end{document}
%====================================================================================================
